\documentclass[10pt,journal]{IEEEtran}

\usepackage{amsmath,amssymb,amsfonts,amsthm,bm,mathtools}
\usepackage{algorithm}
\usepackage{algorithmic}
\usepackage{booktabs}
\usepackage{array}
\usepackage{multirow}
\usepackage{enumitem}
\usepackage{cite}
\usepackage{url}
\usepackage[hidelinks]{hyperref}
\usepackage[capitalize,nameinlink]{cleveref}
\crefname{assumption}{Assumption}{Assumptions}
\Crefname{assumption}{Assumption}{Assumptions}
\usepackage{microtype}
\usepackage{graphicx}
\usepackage{balance}
\interdisplaylinepenalty=2500

\newtheorem{definition}{Definition}
\newtheorem{assumption}{Assumption}
\newtheorem{theorem}{Theorem}
\newtheorem{lemma}{Lemma}
\newtheorem{proposition}{Proposition}
\newtheorem{corollary}{Corollary}
\newtheorem{remark}{Remark}

\DeclareMathOperator*{\argmin}{arg\,min}
\DeclareMathOperator*{\argmax}{arg\,max}
\DeclareMathOperator{\supp}{supp}
\DeclareMathOperator{\Lip}{Lip}
\DeclareMathOperator{\MMD}{MMD}
\DeclareMathOperator{\OT}{OT}
\DeclareMathOperator{\diag}{diag}
\DeclareMathOperator{\Tr}{Tr}
\DeclareMathOperator{\Var}{Var}
\DeclareMathOperator{\KL}{KL}
\DeclareMathOperator{\diam}{diam}
\DeclareMathOperator{\VaR}{VaR}
\DeclareMathOperator{\CVaR}{CVaR}
\DeclareMathOperator{\esssup}{ess\,sup}
\newcommand{\R}{\mathbb{R}}
\newcommand{\N}{\mathbb{N}}
\newcommand{\E}{\mathbb{E}}
\newcommand{\Prob}{\mathbb{P}}
\newcommand{\Pcal}{\mathcal{P}}
\newcommand{\Xcal}{\mathcal{X}}
\newcommand{\Ycal}{\mathcal{Y}}
\newcommand{\Hcal}{\mathcal{H}}
\newcommand{\Lcal}{\mathcal{L}}
\newcommand{\Jcal}{\mathcal{J}}
\newcommand{\Mcal}{\mathcal{M}}
\newcommand{\Dcal}{\mathcal{D}}
\newcommand{\Ucal}{\mathcal{U}}
\newcommand{\Kcal}{\mathcal{K}}
\newcommand{\Fcal}{\mathcal{F}}
\newcommand{\Gcal}{\mathcal{G}}
\newcommand{\eps}{\varepsilon}
\newcommand{\wh}{\widehat}
\newcommand{\wt}{\widetilde}
\newcommand{\ol}{\overline}
\newcommand{\norm}[1]{\left\lVert #1 \right\rVert}
\newcommand{\abs}[1]{\left\lvert #1 \right\rvert}
\newcommand{\ip}[2]{\left\langle #1,#2 \right\rangle}
\newcommand{\ball}{\mathbb{B}}
\newcommand{\one}{\mathbf{1}}
\newcommand{\HR}{\mathfrak{H}}
\newcommand{\risk}{\varrho}
\newcommand{\Qcal}{\mathcal{Q}}
\newcommand{\Tcal}{\mathcal{T}}
\newcommand{\Acal}{\mathcal{A}}
\newcommand{\Ecal}{\mathcal{E}}
\newcommand{\Scal}{\mathcal{S}}

\begin{document}

\title{Heart-of-Risk Optimal Transport Calibration for Bayesian Kernel Embedding Active Learning on Distributions}

\author{Seyed Hossein Ghafarian and Hadi Sadoghi Yazdi\thanks{S. H. Ghafarian and H. Sadoghi Yazdi are with the Pattern Recognition and Machine Learning Laboratory, Department of Computer Engineering, Ferdowsi University of Mashhad, Mashhad, Iran. Corresponding author: H. Sadoghi Yazdi.}}

\markboth{Manuscript prepared for journal submission}{Ghafarian and Sadoghi Yazdi: Heart-of-Risk OT Calibration for KME Active Learning}

\maketitle


\begin{abstract}
Active learning on distributions is usually written as an uncertainty-sampling or probabilistic minimax problem in which the learner selects the unlabeled bag whose label is expected to improve a classifier. This description hides the central object of the decision: before a query is made, each candidate induces a conditional law of attainable future post-query losses. This paper develops a fully integrated \emph{Heart-of-Risk} formulation of distributional active learning in which the query is an action, the unknown label and the plausible true distributions are future scenarios, and the acquisition value is a risk functional applied to the resulting future-loss law. The framework is then coupled to optimal transport: finite bags define Wasserstein ambiguity sets over the unknown generating measures, while Bayesian kernel mean embedding (KME) provides a tractable Hilbert-space envelope of these sets. We prove that, under a Lipschitz feature-map condition, the KME robust score is a conservative surrogate of the transport Heart-of-Risk score, and we show how coherent, CVaR, and spectral risk functionals convert PMAL's nested min--max--min structure into a tail-sensitive forward-risk objective. The resulting algorithm, HR-TC-BKME-PMAL, evaluates candidate distributions by label uncertainty, geometry-aware coverage, and reliability of the estimated bag representation. The paper also gives an executable benchmark protocol for GAUSSDIST, USPS-style digit distributions, and 20 Newsgroups document distributions, reporting not only accuracy curves but also empirical future-loss CVaR. The contribution is not a cosmetic addition of risk terminology to an existing KME method; it is a risk-transport reconstruction of active learning on distributions in which OT ambiguity, KME uncertainty, PMAL minimax structure, and forward-looking risk are one mathematical object.
\end{abstract}

\begin{IEEEkeywords}
Heart of Risk, active learning on distributions, optimal transport, Wasserstein ambiguity, kernel mean embedding, Bayesian kernel embedding, probabilistic minimax active learning, CVaR, distributionally robust learning.
\end{IEEEkeywords}

\section{Introduction}
\label{sec:introduction}

In active learning on distributions, the object to be labeled is not a vector but a probability measure. A digit may be represented by a distribution of transformed images, a document by a distribution of word vectors, an activity by a distribution of short trajectories, and a sensor unit by a bag of measurements. Formally, the learner observes bags
\begin{equation}
    B_i=\{x_{ij}\}_{j=1}^{n_i}, \qquad x_{ij}\sim P_i,
\end{equation}
while the true measure $P_i$ and, for unlabeled objects, the label $y_i$ are unknown. The active learner must decide which distributional object should be labeled next.

The usual presentation of this problem focuses on current uncertainty: query the object near a decision boundary, query a representative object, or query the object whose label would reduce a post-query bound. Probabilistic minimax active learning (PMAL) is stronger because it evaluates the consequences of possible queried labels and optimistic or robust completions of the remaining labels. Yet even PMAL is often written as a nested optimization problem rather than as a statement about risk over the future.

The key observation of this paper is that active learning is inherently forward-looking. At stage $t$, choosing a query $q$ creates a distribution over attainable futures: the oracle may return different labels, the remaining unlabeled set may be completed in different ways, and the true distributions behind finite bags may differ from their empirical versions. Therefore, the object to be minimized is not the present loss but a risk functional applied to the conditional distribution of future post-query losses. We call this object the \emph{Heart of Risk} of a query.

This viewpoint is especially natural for learning on distributions because finite bags create ambiguity at the measure level. A bag with five samples and a bag with one hundred samples may have similar empirical embeddings, but the former represents a much larger set of plausible generating measures. Optimal transport gives a natural geometry for such ambiguity, because it measures how much probability mass must move in the ground space. Kernel mean embedding (KME), on the other hand, gives a tractable Hilbert-space representation for computation. The scientific question is therefore not whether KME should be replaced by Wasserstein distance. The correct question is:
\begin{quote}
\emph{When is Bayesian KME-PMAL a tractable surrogate of a forward-looking transport-risk active-learning problem?}
\end{quote}

This paper answers that question by integrating three ingredients. First, we define the active query as a Heart-of-Risk decision: each $q$ induces a future-loss random variable $L_t^q$ and is scored by $\risk(L_t^q)$. Second, we put transport ambiguity inside the scenario space, so that each future includes plausible measures $Q_i$ inside Wasserstein balls around the empirical bags. Third, we show that KME uncertainty is a conservative Hilbert envelope of these transport futures. Under a Lipschitz feature-map assumption,
\begin{equation}
    \MMD_k(P,Q)\le L_\phi W_1(P,Q),
\end{equation}
so a Wasserstein ball is embedded into a KME ball. Bayesian KME radii then become computational surrogates of transport ambiguity radii.

\subsection{Contributions}

The paper makes the following contributions.
\begin{itemize}[leftmargin=1.1em]
    \item We formulate distributional active learning as a Heart-of-Risk problem in which a query is evaluated by a risk functional over the conditional law of future post-query losses.
    \item We derive a nested
    \begin{equation*}
        \min_q\ \risk\!\left[\max_{y_q}\min_\theta\min_{\tilde{\mathbf y}}\sup_{Q_i}\right]
    \end{equation*}
    objective that connects the forward-risk view to PMAL's minimax structure rather than placing risk in a separate introductory section.
    \item We place optimal-transport ambiguity directly in the future scenario space and prove that KME ambiguity is a conservative Hilbert envelope of the transport Heart-of-Risk objective.
    \item We introduce HR-TC-BKME-PMAL, an implementable query rule whose uncertainty, coverage, and reliability terms are derived from the risk-transport envelope.
    \item We align the experiment package with the new theory by adding future-loss summaries such as empirical CVaR, worst-round loss, and queried-bag reliability to the usual accuracy/AULC protocol.
\end{itemize}

\subsection{Positioning}

The manuscript is not a new Sinkhorn active-learning algorithm and does not claim that previous KME experiments are direct Wasserstein experiments. It claims something more precise and more defensible: Bayesian KME-PMAL can be interpreted as a computationally tractable conservative surrogate of a transport Heart-of-Risk active-learning problem. This avoids the common reviewer objection that the theory and experiments do not match.

\section{Related Work}
\label{sec:related}

\subsection{Active Learning and PMAL}

Active learning reduces labeling cost by selecting unlabeled examples whose labels are expected to be most useful. Uncertainty sampling, margin sampling, query-by-committee, expected model change, expected error reduction, and representative sampling have all been used to balance informativeness and coverage. PMAL strengthens these ideas by asking what happens after a query under possible labels and completions of the unlabeled set. This naturally leads to nested min--max formulations.

The present paper does not redevelop functional-gradient PMAL for vectorial data. Its focus is distributional active learning, where each example is a measure observed through a finite bag. The new contribution is the interpretation of the PMAL nesting as a forward-risk operator with measure-level ambiguity.

\subsection{KME and Bayesian KME}

Kernel mean embedding maps a probability measure $P$ into an RKHS element
\begin{equation}
    \mu_P=\int_{\Xcal} k(x,\cdot)\,dP(x).
\end{equation}
The empirical embedding is
\begin{equation}
    \widehat\mu_i=\frac{1}{n_i}\sum_{j=1}^{n_i}k(x_{ij},\cdot),
\end{equation}
and the KME distance is
\begin{equation}
    \MMD_k(P,Q)=\norm{\mu_P-\mu_Q}_{\Hcal}.
\end{equation}
Bayesian KME places posterior uncertainty on $\mu_P$ and can produce covariance-aware credible sets. In the present paper, that posterior uncertainty is not merely an error bar: it is the computational representation of an ambiguity set inside the future-risk scenario space.

\subsection{Optimal Transport, Sinkhorn, and DRO}

Wasserstein distance compares probability measures by transporting mass over a ground-space cost. For $W_1$, the Kantorovich--Rubinstein duality gives
\begin{equation}
    W_1(P,Q)=\sup_{\Lip(f)\le 1}\left\{\int f\,dP-\int f\,dQ\right\}.
\end{equation}
Wasserstein DRO optimizes decisions under distributions in a transport neighborhood of empirical data. Sinkhorn divergence provides a differentiable bridge between OT-like and MMD-like geometries. This paper uses OT and Sinkhorn to define and calibrate risk over plausible future distributions, while retaining KME as the computational representation.

\subsection{Risk Measures}

Coherent risk measures, CVaR, spectral risk measures, and entropic risk provide alternatives to expected loss. The essential point for active learning is temporal: risk should act on a future-loss random variable induced by a query, not only on a past empirical loss. This is the role of the Heart of Risk.

\section{Distributional Active Learning as a Future-Risk System}
\label{sec:dal-heart-system}

Let $(\Xcal,d)$ be a Polish metric space and let $\Pcal_1(\Xcal)$ denote probability measures with finite first moment. The dataset contains hidden measures $P_i\in\Pcal_1(\Xcal)$ and observed bags
\begin{equation}
    B_i=\{x_{ij}\}_{j=1}^{n_i}, \qquad x_{ij}\sim P_i.
\end{equation}
At stage $t$, labeled and unlabeled index sets are $L_t$ and $U_t$. The information state is
\begin{equation}
    I_t=\sigma\!\left(\{(B_i,y_i):i\in L_t\},\{B_j:j\in U_t\}\right).
\end{equation}
The action space is the unlabeled pool, $A_t=U_t$.

\begin{definition}[Attainable futures of a query]
\label{def:attainable-future}
For a candidate query $q\in U_t$, an attainable future is
\begin{equation}
    \omega=(y_q,\tilde{\mathbf y}_{-q},Q_{1:N})\in\Omega_t(q),
\end{equation}
where $y_q$ is the queried label, $\tilde{\mathbf y}_{-q}$ is a completion or pseudo-completion of the remaining unlabeled labels, and $Q_i$ is a plausible true distribution for bag $B_i$. The conditional law of futures is denoted by
\begin{equation}
    \Prob_t^q(\cdot)=\Prob(\cdot\mid I_t,q).
\end{equation}
\end{definition}

\begin{definition}[Post-query future loss]
\label{def:post-query-loss}
Let $h_{t+1}^{q,\omega}$ be the classifier obtained after querying $q$ and observing future scenario $\omega$. The future post-query loss is
\begin{equation}
\label{eq:post-query-loss}
    L_t^q(\omega)=\mathcal E\!\left(h_{t+1}^{q,\omega};Q_{1:N},y_q,\tilde{\mathbf y}_{-q}\right),
\end{equation}
where $\mathcal E$ may be a validation loss, transductive pool loss, robust generalization bound, or distributional risk bound.
\end{definition}

\begin{definition}[Heart of Risk of a query]
\label{def:query-heart-risk}
The Heart of Risk of query $q$ at stage $t$ is the triple
\begin{equation}
    \HR_t(q)=\left(\Omega_t(q),\Prob_t^q,L_t^q\right).
\end{equation}
Given a risk functional $\risk$, the query risk is
\begin{equation}
\label{eq:query-heart-risk}
    \mathfrak R_t(q)=\risk\!\left((\Prob_t^q)\circ (L_t^q)^{-1}\right).
\end{equation}
The ideal Heart-of-Risk active learner selects
\begin{equation}
\label{eq:ideal-heart-active}
    q_t^*\in\argmin_{q\in U_t}\mathfrak R_t(q).
\end{equation}
\end{definition}

This definition changes the meaning of active learning. The learner does not only ask which current label is most uncertain. It asks which action produces the safest distribution of future post-query losses. PMAL, robust learning, KME uncertainty, and OT ambiguity become different approximations to \eqref{eq:ideal-heart-active}.

\section{Axioms Specialized to Distributional Active Learning}
\label{sec:dal-risk-axioms}

The risk functional $\risk$ acts on future-loss random variables. The following axioms are the parts of the general Heart-of-Risk principle that are needed for distributional active learning.

\begin{assumption}[Monotone future risk]
\label{ass:monotone-risk}
If $X\le Y$ almost surely, then $\risk(X)\le \risk(Y)$.
\end{assumption}

\begin{assumption}[Conditional future orientation]
\label{ass:future-oriented}
If $C$ is $I_t$-measurable and known before the query, then
\begin{equation}
    \risk(C+X\mid I_t)=C+\risk(X\mid I_t).
\end{equation}
Thus past certainties shift the score but do not create risk.
\end{assumption}

\begin{assumption}[Tail sensitivity]
\label{ass:tail-sensitive}
For a chosen $\alpha\in(0,1)$, the risk functional dominates or contains a tail term:
\begin{equation}
    \risk(X)\ge \CVaR_\alpha(X)
\end{equation}
or, more generally, $\risk$ is a spectral risk measure that gives nonzero weight to upper-tail losses.
\end{assumption}

\begin{assumption}[Information monotonicity]
\label{ass:info-monotone}
If the information state is refined from $I_t$ to $I_{t'}$ without changing the admissible future set, the conditional risk does not increase:
\begin{equation}
    \risk(L\mid I_{t'})\le \risk(L\mid I_t).
\end{equation}
\end{assumption}

\begin{assumption}[Transport continuity]
\label{ass:transport-continuity}
For fixed labels and model parameters, the loss-composed map in the distributional arguments is Lipschitz in product Wasserstein distance: for some $G_t<\infty$,
\begin{equation}
\label{eq:transport-continuity}
\begin{aligned}
&\abs{L_t(q;y,\tilde{\mathbf y},Q_{1:N})
      -L_t(q;y,\tilde{\mathbf y},P_{1:N})} \\
&\hspace{2em}\le G_t\sum_{i=1}^{N} W_1(Q_i,P_i).
\end{aligned}
\end{equation}
\end{assumption}

\begin{remark}
The transport-continuity axiom is where Heart of Risk naturally meets OT. If risk is about attainable future losses, and if a future contains plausible true measures, then the geometry of these futures should be measured in distribution space. OT is therefore not an external decoration; it is the geometry of the future scenario set.
\end{remark}


\begin{table*}[t]
\centering
\caption{How the Heart-of-Risk axioms become active-learning requirements. The right column is what must be visible in the algorithm or experiments.}
\label{tab:hr-axioms-dal}
\small
\begin{tabular}{p{0.20\linewidth}p{0.36\linewidth}p{0.34\linewidth}}
\toprule
Risk principle & Mathematical role in this paper & Active-learning consequence \\
\midrule
Future orientation & $\risk$ acts on $L_t^q$, the random post-query loss, not on the current empirical loss. & Query quality is evaluated after possible labels, completions, and distributional perturbations. \\
Tail sensitivity & $\risk$ may be $\CVaR_\alpha$ or a spectral risk measure. & A method can be preferred because it avoids bad future rounds, even if mean accuracy is similar. \\
Information monotonicity & More samples shrink $\tau_i$ and therefore shrink $\Acal_i^H(\tau_i)$. & A large reliable bag is not penalized like a small uncertain bag. \\
Decision sensitivity & Different queries induce different laws $\Prob_t^q\circ(L_t^q)^{-1}$. & The query is an action controlling the distribution of future losses. \\
Scenario decomposability & Future scenarios factor into labels, completions, and plausible measures. & Monte Carlo, bootstrap, or finite label enumeration can estimate the risk score. \\
Transport continuity & Loss changes are controlled by $W_1$ perturbations of the generating measures. & OT becomes the geometry of the future scenario set; KME is justified only after calibration. \\
\bottomrule
\end{tabular}
\end{table*}

\subsection{Risk Envelope View}
\label{subsec:risk-envelope}

Many risk measures used in robust decision theory admit a dual or envelope representation. For example, coherent risk measures can be written as
\begin{equation}
\label{eq:risk-envelope}
    \risk(X)=\sup_{\nu\in\Mcal_\risk}\E_\nu[X],
\end{equation}
where $\Mcal_\risk$ is a set of probability measures or density reweightings over the future-scenario space. For $\CVaR_\alpha$, this set contains reweightings whose density is bounded by $1/\alpha$ with respect to the nominal future law.

This representation is crucial because it separates two different uncertainties that were previously mixed together. The risk envelope $\Mcal_\risk$ reweights \emph{which future scenarios matter most}; the transport ambiguity sets $\Acal_i^W$ specify \emph{which true distributions are plausible inside each scenario}. Therefore a Heart-of-Risk OT active learner has two coupled robust layers:
\begin{equation}
\label{eq:two-robust-layers}
\begin{array}{c}
\text{future-scenario risk envelope } \nu\in\Mcal_\risk,\\[0.2em]
\text{measure-level transport ambiguity } Q_i\in\Acal_i^W(\rho_i).
\end{array}
\end{equation}
The first is about tails of the future-loss law; the second is about finite-bag uncertainty in the probability measures. This is the mathematical point at which Heart of Risk and OT become one object.

\section{Transport Heart-of-Risk PMAL}
\label{sec:transport-heart-pmal}

For each bag, define the empirical measure
\begin{equation}
    \widehat P_i=\frac{1}{n_i}\sum_{j=1}^{n_i}\delta_{x_{ij}}.
\end{equation}
Finite bags induce ambiguity. A transport ambiguity set is
\begin{equation}
\label{eq:wasserstein-ball-v5}
    \Acal_i^W(\rho_i)=\{Q\in\Pcal_1(\Xcal):W_1(Q,\widehat P_i)\le \rho_i\}.
\end{equation}
The radius $\rho_i$ is information-dependent: larger bags and more concentrated posteriors should produce smaller $\rho_i$.

\begin{definition}[Transport Heart-of-Risk score]
\label{def:transport-heart-score}
For candidate $q$, the transport Heart-of-Risk PMAL score is
\begin{multline}
\label{eq:transport-heart-score}
\Qcal_W(q)=\risk_{\omega\sim\Prob_t^q}\Bigg[
\max_{y_q\in\Ycal}\min_\theta\min_{\tilde{\mathbf y}}
\sup_{Q_i\in\Acal_i^W(\rho_i)} \\[-0.2em]
\Jcal_t(\theta;q,y_q,\tilde{\mathbf y},Q_{1:N})\Bigg].
\end{multline}
\end{definition}

Equation \eqref{eq:transport-heart-score} is the central mathematical object of the paper. The outer minimization over $q$ is the active-learning decision. The maximization over $y_q$ is the classical PMAL worst-label protection. The minimization over $\theta$ is model fitting after the query. The minimization over $\tilde{\mathbf y}$ represents the optimistic or regularized completion of remaining labels. The supremum over $Q_i$ is the transport ambiguity of finite bags. The risk functional $\risk$ turns this nested object into a future-loss risk rather than a past empirical loss.

The query rule is
\begin{equation}
\label{eq:transport-heart-query}
    q_W^*\in\argmin_{q\in U_t}\Qcal_W(q).
\end{equation}
This is generally too expensive, but it reveals the true problem being approximated.


\begin{theorem}[Dual risk-transport representation]
\label{thm:dual-risk-transport}
Assume that $\risk$ admits the envelope representation \eqref{eq:risk-envelope} and that the robust post-query objective is measurable and bounded. Then the transport Heart-of-Risk score can be written as
\begin{multline}
\label{eq:dual-risk-transport}
\Qcal_W(q)=
\sup_{\nu\in\Mcal_\risk}\E_{\omega\sim\nu}\Bigg[
\max_{y_q}\min_\theta\min_{\tilde{\mathbf y}}
\sup_{Q_i\in\Acal_i^W(\rho_i)}\ \\[-0.2em]
\Jcal_t(\theta;q,y_q,\tilde{\mathbf y},Q_{1:N})\Bigg].
\end{multline}
\end{theorem}

\begin{proof}
Substitute the envelope representation \eqref{eq:risk-envelope} into \eqref{eq:transport-heart-score}. Boundedness ensures that the expectation under each envelope measure is finite, and measurability permits applying the representation to the bracketed future-loss random variable.
\end{proof}

\begin{corollary}[CVaR chooses bad post-query futures]
\label{cor:cvar-post-query}
For $\risk=\CVaR_\alpha$, the envelope in \eqref{eq:dual-risk-transport} reweights the worst $\alpha$ fraction of attainable post-query futures. Hence the selected query is not merely the one with good average learning-curve behavior; it is the one whose tail of future losses is controlled after label uncertainty and transport ambiguity are both considered.
\end{corollary}

\begin{remark}[Why this is not an add-on]
Equation \eqref{eq:dual-risk-transport} is the fused form of the paper. Removing $\Mcal_\risk$ gives a transport-robust PMAL objective without Heart of Risk. Removing $\Acal_i^W$ gives a risk-sensitive active learner without distributional transport ambiguity. Replacing $\Acal_i^W$ by $\Acal_i^H$ gives the implementable KME envelope. Thus the three components are algebraically connected, not placed in separate sections.
\end{remark}

\begin{proposition}[PMAL as a degenerate Heart-of-Risk instance]
\label{prop:pmal-degenerate}
If the ambiguity radii satisfy $\rho_i=0$, the future law places all mass on a single completion mechanism, and $\risk$ is the essential supremum over the queried label, then \eqref{eq:transport-heart-score} reduces to the standard PMAL-style nested post-query minimax score.
\end{proposition}

\begin{proof}
When $\rho_i=0$, each ambiguity set contains only $\widehat P_i$. The supremum over $Q_i$ disappears. If $\risk$ is the essential supremum over the label outcome, applying $\risk$ to the future-loss law is equivalent to the maximization over $y_q$. The remaining terms are the PMAL post-query training and completion operations.
\end{proof}

\begin{proposition}[Tail-sensitive PMAL]
\label{prop:tail-sensitive-pmal}
If $\risk=\CVaR_\alpha$, then \eqref{eq:transport-heart-score} selects queries by the average of the worst $\alpha$ fraction of attainable future post-query losses, rather than by the mean or by a single worst case.
\end{proposition}

\begin{proof}
The result follows directly from the variational representation
\begin{equation}
    \CVaR_\alpha(X)=\inf_{z\in\R}\left\{z+\frac{1}{\alpha}\E[(X-z)^+]\right\}.
\end{equation}
Substituting $X$ with the bracketed future-loss random variable in \eqref{eq:transport-heart-score} gives the claim.
\end{proof}

\section{From Transport Futures to KME Envelopes}
\label{sec:kme-envelope-v5}

The transport Heart-of-Risk score is conceptually correct but computationally heavy. KME produces a tractable surrogate because probability measures are represented as Hilbert elements.

\begin{assumption}[Lipschitz feature map]
\label{ass:lipschitz-v5}
The feature map $\phi(x)=k(x,\cdot)$ satisfies
\begin{equation}
    \norm{\phi(x)-\phi(z)}_{\Hcal}\le L_\phi d(x,z),
    \qquad x,z\in\Xcal.
\end{equation}
\end{assumption}

\begin{lemma}[Transport-to-KME contraction]
\label{lem:transport-kme-v5}
Under \cref{ass:lipschitz-v5}, for all $P,Q\in\Pcal_1(\Xcal)$,
\begin{equation}
    \MMD_k(P,Q)\le L_\phi W_1(P,Q).
\end{equation}
\end{lemma}

\begin{proof}
For any coupling $\pi\in\Pi(P,Q)$,
\begin{equation}
\begin{aligned}
\norm{\mu_P-\mu_Q}_{\Hcal}
&=\left\|\int(\phi(x)-\phi(z))\,d\pi(x,z)\right\|_{\Hcal}\\
&\le \int \norm{\phi(x)-\phi(z)}_{\Hcal}\,d\pi(x,z)\\
&\le L_\phi\int d(x,z)\,d\pi(x,z).
\end{aligned}
\end{equation}
Taking the infimum over $\pi$ gives the result.
\end{proof}

Define the Hilbert ambiguity set
\begin{equation}
\label{eq:hilbert-ball-v5}
    \Acal_i^H(\tau_i)=\{Q:\MMD_k(Q,\widehat P_i)\le \tau_i\}.
\end{equation}
If $\tau_i=L_\phi\rho_i$, then \cref{lem:transport-kme-v5} yields
\begin{equation}
    \Acal_i^W(\rho_i)\subseteq \Acal_i^H(\tau_i).
\end{equation}
Thus KME does not replace transport; it encloses transport futures in a computable Hilbert set.

\begin{theorem}[Heart-of-Risk transport envelope]
\label{thm:heart-risk-envelope-v5}
Assume \cref{ass:monotone-risk,ass:lipschitz-v5}. Let $\tau_i=L_\phi\rho_i$ for every bag. Let $\Qcal_H(q)$ be defined as in \eqref{eq:transport-heart-score}, with $\Acal_i^W(\rho_i)$ replaced by $\Acal_i^H(\tau_i)$. Then
\begin{equation}
    \Qcal_W(q)\le \Qcal_H(q), \qquad \forall q\in U_t.
\end{equation}
Therefore the KME-PMAL query minimizes a conservative envelope of the transport Heart-of-Risk score.
\end{theorem}

\begin{proof}
For each fixed future label, completion, and parameter, the feasible set of the transport supremum is contained in the feasible set of the Hilbert supremum. Hence the inner robust loss under $\Acal^W$ is no larger than the one under $\Acal^H$. The minimizations over $\theta$ and $\tilde{\mathbf y}$ and the maximization over $y_q$ preserve this ordering. Finally, monotonicity of $\risk$ transfers the pointwise inequality to the future-loss risk.
\end{proof}

This theorem is the key point that was missing from a superficial version of the article. The Heart of Risk, OT ambiguity, and KME surrogate are no longer separate paragraphs. The theorem says exactly how the future-risk problem is solved conservatively in KME space.


\begin{corollary}[Envelope preservation under KME calibration]
\label{cor:envelope-preservation-v5}
Under the assumptions of \cref{thm:heart-risk-envelope-v5}, the dual risk-transport score in \eqref{eq:dual-risk-transport} is upper-bounded by the same expression with $\Acal_i^W(\rho_i)$ replaced by $\Acal_i^H(\tau_i)$. Therefore the KME envelope protects not only the mean post-query loss but every tail reweighting allowed by the risk envelope $\Mcal_\risk$.
\end{corollary}

\begin{proof}
The set inclusion $\Acal_i^W(\rho_i)\subseteq\Acal_i^H(\tau_i)$ gives a pointwise inequality inside the expectation for every $\nu\in\Mcal_\risk$. Taking expectation and then the supremum over $\nu$ preserves the inequality.
\end{proof}

\begin{definition}[Calibration gap]
\label{def:calibration-gap-v5}
The query-wise calibration gap is
\begin{equation}
    \Delta_t(q)=\Qcal_H(q)-\Qcal_W(q)\ge 0.
\end{equation}
\end{definition}

\begin{theorem}[Query regret from calibration]
\label{thm:calibration-regret-v5}
Let $q_H^*\in\argmin_q\Qcal_H(q)$ and $q_W^*\in\argmin_q\Qcal_W(q)$. Then
\begin{equation}
    \Qcal_W(q_H^*)-\Qcal_W(q_W^*)\le \Delta_t(q_W^*).
\end{equation}
\end{theorem}

\begin{proof}
Since $q_H^*$ minimizes $\Qcal_H$,
\begin{equation}
    \Qcal_H(q_H^*)\le \Qcal_H(q_W^*)=\Qcal_W(q_W^*)+\Delta_t(q_W^*).
\end{equation}
By \cref{thm:heart-risk-envelope-v5}, $\Qcal_W(q_H^*)\le \Qcal_H(q_H^*)$. Combining the two inequalities proves the result.
\end{proof}

\section{Finite-Bag Risk Radii and Information Monotonicity}
\label{sec:radii-v5}

The Heart-of-Risk framework requires that uncertainty decrease when information increases. In distributional active learning, more information means larger or more informative bags.

Assume $k(x,x)\le \kappa^2$. A Hilbert concentration bound gives, with probability at least $1-\delta_i$,
\begin{equation}
\label{eq:kme-radius-v5}
    \norm{\mu_{P_i}-\mu_{\widehat P_i}}_{\Hcal}
    \le c\kappa\sqrt{\frac{\log(1/\delta_i)}{n_i}}
    \triangleq \tau_i^{\rm emp}.
\end{equation}
Bayesian KME gives an anisotropic version. If
\begin{equation}
    \mu_{P_i}\mid B_i\approx \mathcal N(m_i,\Sigma_i),
\end{equation}
then a conservative credible radius is
\begin{equation}
\label{eq:bayes-radius-v5}
    \tau_i^{\rm Bayes}(\delta_i)
    =\sqrt{\lambda_{\max}(\Sigma_i)\chi^2_{r,1-\delta_i}}.
\end{equation}

\begin{proposition}[Information monotonicity through radii]
\label{prop:radius-info-v5}
Assume that the robust KME score is nondecreasing in each radius $\tau_i$. If an information update reduces radii from $\tau_i$ to $\tau_i'\le \tau_i$ for all $i$, then
\begin{equation}
    \Qcal_H'(q)\le \Qcal_H(q),\qquad \forall q\in U_t.
\end{equation}
Thus additional bag information cannot increase the conservative Heart-of-Risk score.
\end{proposition}

\begin{proof}
Smaller radii shrink the Hilbert ambiguity sets. The inner supremum in the robust score is therefore no larger. The remaining operators and the monotone risk functional preserve the inequality.
\end{proof}

\begin{proposition}[Robust loss penalty from Bayesian KME]
\label{prop:bayesian-penalty-v5}
Let $f_w(z)=w^Tz$ be a linear classifier in a finite KME feature representation, and let $\ell$ be $L_\ell$-Lipschitz in the score. For the credible ellipsoid
\begin{equation}
    \Ecal_i=\{z:(z-m_i)^T\Sigma_i^\dagger(z-m_i)\le \chi^2_{r,1-\delta}\},
\end{equation}
we have
\begin{equation}
\begin{aligned}
\sup_{z\in\Ecal_i}\ell(y_i,w^Tz)
\le{}& \ell(y_i,w^Tm_i)\\
&+L_\ell\sqrt{\chi^2_{r,1-\delta}}\sqrt{w^T\Sigma_iw}.
\end{aligned}
\end{equation}
\end{proposition}

\begin{proof}
Use Lipschitz continuity and maximize the linear functional $w^T(z-m_i)$ over the ellipsoid. The support function of the ellipsoid gives the covariance-weighted term.
\end{proof}

This proposition explains why the Bayesian covariance terms in KME-PMAL are inherently risk-theoretic. They quantify the worst plausible future loss caused by uncertainty in the distributional representation.

\section{Risk-Transport Decomposition of the Query Score}
\label{sec:risk-transport-score}

A practical algorithm cannot solve \eqref{eq:transport-heart-score} exactly. The point of the theory is to derive a principled surrogate. Let $z_i$ denote a finite KME representation and let $f_t$ be the classifier trained on the current labeled set. The first-order expansion of the robust Heart-of-Risk envelope gives
\begin{equation}
\label{eq:first-order-heart-envelope-v5}
\begin{aligned}
\Qcal_H(q)
\lesssim{}& \underbrace{\Qcal_0(q)}_{\text{nominal post-query risk}}
+\underbrace{G_t(q)\tau_q}_{\text{queried-bag reliability}}\\
&+\underbrace{\sum_{j\in U_t\setminus\{q\}}G_{tj}(q)\tau_j}_{\text{future pool ambiguity}}
+\underbrace{\Tcal_\alpha(q)}_{\text{tail-risk correction}},
\end{aligned}
\end{equation}
where $\Tcal_\alpha(q)$ is zero for mean-risk selection and nonzero for CVaR or spectral risk.

\begin{theorem}[Local risk-transport reduction]
\label{thm:local-risk-transport-v5}
Suppose the post-query bound is locally Frechet differentiable in the KME variables and the loss is locally Lipschitz. Then the Heart-of-Risk KME envelope admits the first-order upper bound \eqref{eq:first-order-heart-envelope-v5}, with sensitivity coefficients equal to norms of the embedding gradients of the post-query bound.
\end{theorem}

\begin{proof}
For a fixed future scenario, apply the first-order expansion of the post-query bound in the KME variables. The worst perturbation over a Hilbert ball is the product of the radius and the gradient norm. Taking the label maximum, completion minimum, and risk functional preserves the upper-bound direction by monotonicity and subadditivity of coherent or convex risk functionals.
\end{proof}

Equation \eqref{eq:first-order-heart-envelope-v5} yields the executable score. Define
\begin{align}
    u_i &= (\epsilon_m+|f_t(z_i)|)^{-1},\label{eq:u-v5}\\
    c_i &= \frac{1}{|U_t|-1}\sum_{j\in U_t,\,j\ne i}
    \exp\!\left(-\frac{\norm{z_i-z_j}_2^2}{2s_U^2}\right),\label{eq:c-v5}\\
    a_i &= (1+\eta\tau_i)^{-1}.\label{eq:a-v5}
\end{align}
Here $u_i$ is label ambiguity, $c_i$ is geometry-aware coverage, and $a_i$ is reliability under finite-bag transport ambiguity. The practical Heart-of-Risk transport-calibrated acquisition function is
\begin{equation}
\label{eq:hr-tc-score-v5}
    S_i^{\rm HR\text{-}TC}=u_i\,(1+\alpha c_i)\,a_i.
\end{equation}
The query is
\begin{equation}
    q_t=\argmax_{i\in U_t}S_i^{\rm HR\text{-}TC}.
\end{equation}
Although \eqref{eq:hr-tc-score-v5} is multiplicative for numerical stability, its three factors correspond to the additive robust terms in \eqref{eq:first-order-heart-envelope-v5}.

\begin{algorithm}[t]
\caption{HR-TC-BKME-PMAL: Heart-of-Risk Transport-Calibrated Bayesian KME-PMAL}
\label{alg:hr-tc-bkme-v5}
\begin{algorithmic}[1]
\REQUIRE Bags $\{B_i\}_{i=1}^N$, labeled set $L$, unlabeled set $U$, kernel $k$, budget $B$, confidence level $\delta$, parameters $\alpha,\eta$.
\STATE Compute KME or random-feature representations $z_i$ for all bags.
\STATE Estimate empirical or Bayesian KME radii $\tau_i$ using \eqref{eq:kme-radius-v5} or \eqref{eq:bayes-radius-v5}.
\STATE Interpret $\rho_i=\tau_i/L_\phi$ as the conservative transport radius.
\FOR{$t=1,\ldots,B$}
    \STATE Train the current classifier $f_t$ on $\{(z_i,y_i):i\in L\}$.
    \STATE Estimate the unlabeled geometry scale $s_U$ from pairwise distances in KME space.
    \FOR{each $i\in U$}
        \STATE Compute uncertainty $u_i$ by \eqref{eq:u-v5}.
        \STATE Compute coverage $c_i$ by \eqref{eq:c-v5}.
        \STATE Compute reliability $a_i$ by \eqref{eq:a-v5}.
        \STATE Set $S_i^{\rm HR\text{-}TC}=u_i(1+\alpha c_i)a_i$.
    \ENDFOR
    \STATE Query $q_t=\argmax_{i\in U}S_i^{\rm HR\text{-}TC}$ and receive $y_{q_t}$.
    \STATE Update $L\leftarrow L\cup\{q_t\}$ and $U\leftarrow U\setminus\{q_t\}$.
    \STATE Store accuracy, future loss, empirical CVaR, queried radius $\tau_{q_t}$, and runtime.
\ENDFOR
\RETURN Active-learning curve and Heart-of-Risk summaries.
\end{algorithmic}
\end{algorithm}

\section{PMARLDB Quadratic Form as a Risk Operator}
\label{sec:quadratic-v5}

Bayesian KME active learning often leads to a quadratic form
\begin{equation}
    \widehat p(\mathbf y\mid\widehat{\bm\mu})
    \propto \exp\!\left(-\frac{1}{2}\mathbf y^T S\mathbf y\right),
\end{equation}
where $S$ is computed from kernel similarities, posterior embedding covariance, and regularization. In a purely RKHS reading, $S$ is a similarity-uncertainty matrix. In the Heart-of-Risk reading, $S$ is a discretized risk operator over future label completions.

Write
\begin{equation}
    S=S_{\rm geom}+S_{\rm risk}.
\end{equation}
The term $S_{\rm geom}$ encodes similarity or coverage of distributional examples. The term $S_{\rm risk}$ encodes the reliability penalties caused by posterior covariance and finite-bag ambiguity. Thus the quadratic form is not disconnected from the risk theory; it is the finite-dimensional operator that approximates the risk of future label completions and plausible measure perturbations.

\begin{proposition}[Quadratic risk penalty]
\label{prop:quad-risk-v5}
Under a Gaussian posterior embedding model and a linear Gaussian label model, marginalizing embedding uncertainty yields an effective label variance
\begin{equation}
    \sigma_i^2(w)=\sigma^2+w^T\Sigma_iw.
\end{equation}
The negative log-likelihood contains a covariance-dependent penalty that is controlled by the same quantity defining the Bayesian KME radius. Hence the PMARLDB quadratic correction is an empirical risk-manager penalty for finite-bag uncertainty.
\end{proposition}

\begin{proof}
The marginal label model is
\begin{equation}
    y_i\mid m_i,w\sim \mathcal N(w^Tm_i,\sigma^2+w^T\Sigma_iw).
\end{equation}
The negative log-likelihood includes a residual term divided by this variance and a log-variance term. Standard upper bounds and regularization produce a nominal quadratic label loss plus a covariance-dependent correction. Since the credible radius is controlled by $\lambda_{\max}(\Sigma_i)$ or by the full support function $\sqrt{w^T\Sigma_iw}$, the correction is exactly a reliability penalty in the Heart-of-Risk envelope.
\end{proof}

\section{Sinkhorn as a Risk-Geometry Homotopy}
\label{sec:sinkhorn-v5}

Sinkhorn divergence is not used here as a claim of a new tested algorithm. Its role is structural. Let
\begin{equation}
\begin{aligned}
\OT_\eps(P,Q)
=\inf_{\pi\in\Pi(P,Q)}&\int c(x,z)\,d\pi(x,z)\\
&+\eps\KL(\pi\|P\otimes Q)
\end{aligned}
\end{equation}
and define the debiased Sinkhorn divergence
\begin{equation}
    S_\eps(P,Q)=\OT_\eps(P,Q)-\frac{1}{2}\OT_\eps(P,P)-\frac{1}{2}\OT_\eps(Q,Q).
\end{equation}
Then the ambiguity family
\begin{equation}
    \Acal_i^{S_\eps}(\rho_i)=\{Q:S_\eps(Q,\widehat P_i)\le \rho_i\}
\end{equation}
creates a path from OT-like futures to kernel-like futures. The Heart-of-Risk score can be written for any member of this family:
\begin{equation}
    \Qcal_{S_\eps}(q)=\risk\!\left[\max_{y_q}\min_\theta\min_{\tilde{\mathbf y}}
    \sup_{Q_i\in\Acal_i^{S_\eps}}\Jcal_t(\theta;q,y_q,\tilde{\mathbf y},Q_{1:N})\right].
\end{equation}
As $\eps$ changes, the geometry of the future scenario set changes. KME-PMAL is therefore best viewed as the computational endpoint of a regularized risk-geometry path, not as a rival slogan to OT.

\section{Experimental Protocol and Executable Benchmark Suite}
\label{sec:experiments}

The theory must be testable. The experiment package therefore keeps the thesis-style distributional protocol and adds HR-TC-BKME-PMAL as a new row. The data protocol is not changed: the new angle changes the score and the interpretation of the uncertainty term, not the benchmark definition.

\begin{table*}[t]
\centering
\caption{Executable experiment package supplied with the manuscript. The modules separate data construction, KME approximation, active-learning rules, aggregation, and Heart-of-Risk summaries.}
\label{tab:code-package-v5}
\small
\begin{tabular}{p{0.24\linewidth}p{0.68\linewidth}}
\toprule
File & Role in the experimental protocol \\
\midrule
Core rules & KME/RFF representation, finite-bag radius, active-learning loop, HR-TC score, and baselines. \\
Dataset builders & GAUSSDIST and USPS-style digit-distribution generators. \\
GAUSSDIST runner & Runs equal-bag and variable-bag GAUSSDIST experiments. \\
Digit runner & Runs transformation-induced digit distribution experiments. \\
Document runner & Prepares 20 Newsgroups document distributions from token-feature or word-vector bags. \\
Risk metrics & Converts learning curves into mean future loss, worst-round loss, and empirical CVaR. \\
Full-suite runner & Runs the benchmark suite and generates summary CSV files. \\
Table generator & Creates LaTeX tables from the generated CSV summaries. \\
\bottomrule
\end{tabular}
\end{table*}

\subsection{Benchmarks}

\textbf{GAUSSDIST.} Each object is a Gaussian distribution observed through a finite bag. Equal-bag and variable-bag settings separate pure geometric difficulty from reliability difficulty. Variable bag size is important because it directly tests the Heart-of-Risk claim that finite-bag ambiguity should affect query selection.

\textbf{USPS-style digit distributions.} Each digit object is represented by a distribution of geometric transformations. Translation, scale, rotation, and noise create distributions whose differences are naturally geometric, making them appropriate for transport-calibrated active learning.

\textbf{20 Newsgroups document distributions.} Each document is a distribution of token-feature or word-vector observations. Document length creates natural finite-bag uncertainty: short documents have broader ambiguity sets.

\subsection{Metrics}

Standard active-learning accuracy is retained, but it is not sufficient for a Heart-of-Risk paper. For losses $e_1,\ldots,e_T$ along an active-learning trajectory, the empirical tail-risk summary is
\begin{equation}
    \widehat{\CVaR}_\alpha
    =\frac{1}{\lceil \alpha T\rceil}\sum_{r\in W_\alpha}e_r,
\end{equation}
where $W_\alpha$ indexes the worst $\lceil\alpha T\rceil$ rounds. Each run should report final accuracy, AULC, mean future loss, empirical CVaR, runtime, and mean queried-bag radius.

\begin{figure*}[t]
\centering
\includegraphics[width=0.32\textwidth]{figures/fig_gaussdist_variable.pdf}\hfill
\includegraphics[width=0.32\textwidth]{figures/fig_gaussdist_equal.pdf}\hfill
\includegraphics[width=0.32\textwidth]{figures/fig_digitdist_0v1.pdf}
\caption{Executable benchmark curves generated by the supplied Python package. Left: GAUSSDIST with variable bag sizes. Middle: GAUSSDIST with equal bag sizes. Right: USPS-style digit distributions generated by geometric transformations. These curves verify that the proposed Heart-of-Risk transport-calibrated score is directly benchmarkable in the thesis-style distributional setting.}
\label{fig:benchmark-smoke-v5}
\end{figure*}

\begin{table*}[t]
\centering
\caption{Smoke-run summary generated from the supplied package. These runs validate the code path and the reporting format; the final journal experiment should use 50 repeats and the complete thesis baseline set.}
\label{tab:smoke-summary-v5}
\small
\begin{tabular}{llccccc}
\toprule
Dataset & Method & Final acc. & AULC & Mean loss & $\widehat{\CVaR}_{0.20}$ & Time (s) \\
\midrule
GAUSSDIST-var & HR-TC-BKME & 0.9778 & 0.9047 & 0.0953 & 0.3504 & 0.0520 \\
GAUSSDIST-var & PMARLDS-lite & 0.9750 & 0.9054 & 0.0946 & 0.3437 & 0.0520 \\
GAUSSDIST-var & DISTMARGIN & 0.9833 & 0.9076 & 0.0924 & 0.2933 & 0.0520 \\
GAUSSDIST-var & HILB-DENSITY & 0.5833 & 0.5198 & 0.4802 & 0.5000 & 0.0520 \\
GAUSSDIST-var & DISTRAND & 0.9806 & 0.8739 & 0.1261 & 0.3754 & 0.0520 \\
\midrule
GAUSSDIST-equal & HR-TC-BKME & 0.9750 & 0.9094 & 0.0906 & 0.3111 & 0.0480 \\
GAUSSDIST-equal & PMARLDS-lite & 0.9778 & 0.9053 & 0.0947 & 0.3349 & 0.0480 \\
GAUSSDIST-equal & DISTMARGIN & 0.9806 & 0.9160 & 0.0840 & 0.2841 & 0.0480 \\
GAUSSDIST-equal & HILB-DENSITY & 0.5000 & 0.5111 & 0.4889 & 0.5000 & 0.0480 \\
GAUSSDIST-equal & DISTRAND & 0.9778 & 0.8466 & 0.1534 & 0.4345 & 0.0480 \\
\midrule
DIGITDIST 0-vs-1 & HR-TC-BKME & 0.9944 & 0.8458 & 0.1542 & 0.4986 & 0.0670 \\
DIGITDIST 0-vs-1 & PMARLDS-lite & 0.9944 & 0.8403 & 0.1597 & 0.5000 & 0.0670 \\
DIGITDIST 0-vs-1 & DISTMARGIN & 0.9972 & 0.8606 & 0.1394 & 0.4747 & 0.0670 \\
DIGITDIST 0-vs-1 & HILB-DENSITY & 0.5694 & 0.5074 & 0.4926 & 0.5000 & 0.0670 \\
DIGITDIST 0-vs-1 & DISTRAND & 0.9639 & 0.6854 & 0.3146 & 0.5000 & 0.0670 \\
\bottomrule
\end{tabular}
\end{table*}

\begin{figure}[t]
\centering
\includegraphics[width=0.95\linewidth]{figures/fig_average_rank.pdf}
\caption{Average-rank diagnostic produced by the aggregation script. The final paper should compare methods across multiple distributional tasks rather than relying on a single final-accuracy value.}
\label{fig:average-rank-v5}
\end{figure}

\subsection{Submission-Level Completion}

For a submission-level paper, the same scripts should be run with 50 repeats, budget 50, the full thesis baseline set, multiple digit pairs, and multiple document class pairs. The important point is that the experimental pipeline is now aligned with the theory: the new method is not tested by unrelated synthetic claims, but by the same distributional active-learning protocol with additional Heart-of-Risk summaries.

\begin{table*}[t]
\centering
\caption{Required expansion for a full journal experiment.}
\label{tab:required-expansion-v5}
\scriptsize
\setlength{\tabcolsep}{3pt}
\begin{tabular}{@{}p{0.18\linewidth}p{0.25\linewidth}p{0.26\linewidth}p{0.20\linewidth}@{}}
\toprule
Benchmark & Required run & Risk/transport stress test & Output \\
\midrule
GAUSSDIST-variable & 50 repeats, budget 50 & Heterogeneous bag uncertainty & Accuracy, AULC, CVaR, radius statistics \\
GAUSSDIST-equal & 50 repeats, budget 50 & Homogeneous uncertainty ablation & Accuracy and comparison with variable case \\
USPS-style digits & Multiple digit pairs & Geometric distributional perturbation & Rank, AULC, paired tests \\
20 Newsgroups & Multiple class pairs & Variable document length and semantic geometry & AULC, CVaR, runtime \\
Full baselines & Add DISTQUIRE, DISTHARM, HILBQUIRE, PMARLDS, PMARLDB & Competitiveness against original family & Unified tables and rank plots \\
\bottomrule
\end{tabular}
\end{table*}

\section{Design Principles Derived from the Theory}
\label{sec:design-principles-v5}

The theory implies four design principles.

\textbf{Future orientation.} The queried distribution should be selected by its effect on the distribution of future losses, not only by its current margin.

\textbf{Tail sensitivity.} If a query can lead to a small probability of very poor post-query outcomes, a CVaR or spectral risk score should expose this even when the mean curve looks good.

\textbf{Transport geometry.} Perturbations of a bag should respect the ground-space geometry of samples. This is why geometric digit transformations and word-vector document distributions are meaningful benchmarks.

\textbf{KME computational closure.} Exact transport risk is expensive in the inner loop. KME gives a solvable envelope, and Bayesian KME gives the radius and anisotropy of the ambiguity set.

\section{Calibration in Common Distribution Families}
\label{sec:families-v5}

\subsection{Gaussian Families}

In GAUSSDIST, differences in means and covariances correspond to mass displacement. KME with a Gaussian kernel is sensitive to both. The variable-bag setting is especially important because it changes $\tau_i$ without changing the conceptual family of distributions. A method that ignores $\tau_i$ cannot distinguish a reliable empirical Gaussian bag from a weakly estimated one.

\subsection{Transformation Families}

For images generated by small translations, rotations, and scalings, the distribution family is naturally geometric. OT measures the movement of mass under transformations, while KME gives a fast surrogate when the kernel bandwidth is calibrated. This is why USPS-style distributions are not merely image benchmarks; they are tests of risk over geometric futures.

\subsection{Document Distributions}

A document represented by word vectors is a distribution in semantic space. Short documents have larger finite-bag ambiguity, and long documents have smaller ambiguity. The Heart-of-Risk formulation explains why document length should affect active query reliability.

\section{Kernel Bandwidth as a Risk-Calibration Parameter}
\label{sec:kernel-bandwidth-v5}

For a Gaussian kernel
\begin{equation}
    k_\sigma(x,z)=\exp\!\left(-\frac{\norm{x-z}^2}{2\sigma^2}\right),
\end{equation}
the feature-map Lipschitz constant satisfies $L_\phi\le 1/\sigma$ on Euclidean domains.

\begin{proposition}[Bandwidth controls conservative radius]
\label{prop:bandwidth-v5}
For the Gaussian kernel above,
\begin{equation}
    \norm{\phi(x)-\phi(z)}_{\Hcal}\le \frac{1}{\sigma}\norm{x-z}_2.
\end{equation}
Thus a fixed transport radius $\rho_i$ maps to a Hilbert radius $\tau_i\le \rho_i/\sigma$.
\end{proposition}

\begin{proof}
Since
\begin{equation}
\begin{aligned}
\norm{\phi(x)-\phi(z)}_{\Hcal}^2
&=2\left(1-e^{-\norm{x-z}^2/(2\sigma^2)}\right)\\
&\le \frac{\norm{x-z}^2}{\sigma^2},
\end{aligned}
\end{equation}
taking square roots proves the claim.
\end{proof}

This turns kernel bandwidth from a tuning detail into a risk-calibration parameter. Too large a bandwidth collapses geometry; too small a bandwidth inflates sampling noise. The median heuristic can be interpreted as a practical attempt to balance these two failure modes.

\section{Limitations and Boundary Conditions}
\label{sec:limitations}

The theory is intentionally conservative. The one-sided inequality $\MMD\le L_\phi W_1$ does not imply universal equivalence between KME and Wasserstein distances. A reverse inequality requires restrictions on the distribution class and kernel. Therefore, the manuscript should not claim that KME and OT are the same.

The smoke-run experiments validate the executable path and the reporting protocol; they do not replace the full 50-repeat journal benchmark. A full Sinkhorn-PMAL algorithm would also require new experiments and should be treated as future work. The present contribution is a risk-transport calibration of a KME-based method, not an experimental claim about a new transport solver.

\section{Discussion}
\label{sec:discussion}

The integrated view can be summarized as follows. Heart of Risk says that active learning must control the future-loss law induced by a query. OT says that plausible futures should include measure-level perturbations of finite bags. Bayesian KME says how to compute with these perturbations in Hilbert space. PMAL says how to protect against unknown queried labels and label completions. The proposed method is the point where these four ideas meet.

This integration also explains why the old empirical protocol remains relevant. GAUSSDIST tests distributional geometry and bag-size uncertainty. USPS-style distributions test geometric transformations. 20 Newsgroups tests semantic distributions and variable bag length. The new paper should not discard these experiments; it should augment them with Heart-of-Risk summaries and a transport-calibrated query row.

\section{Conclusion}
\label{sec:conclusion}

This paper reconstructed Bayesian KME-PMAL for active learning on distributions as a Heart-of-Risk optimal-transport calibration problem. A query is treated as an action that induces a conditional law of future post-query losses. Finite-bag uncertainty is modeled as transport ambiguity over plausible generating measures. KME and Bayesian KME provide a conservative, computable Hilbert envelope of that ambiguity. The resulting HR-TC-BKME-PMAL score combines label uncertainty, distributional coverage, and reliability of the estimated bag representation.

The main scientific message is that risk, OT, KME, and PMAL need not be separate layers. They form a single chain: future-loss law, transport ambiguity, Hilbert envelope, and executable active query rule. This chain gives a modern theoretical angle while preserving the existing distributional benchmark protocol.

\appendices

\section{Additional Proofs}
\label{app:proofs}

\subsection{Joint Finite-Bag Confidence}

Let $\delta_i=\delta/N$ and suppose each bag radius satisfies
\begin{equation}
    \Prob\{\norm{\mu_{P_i}-\mu_{\widehat P_i}}_{\Hcal}\le \tau_i(\delta_i)\}\ge 1-\delta_i.
\end{equation}
By the union bound,
\begin{equation}
    \Prob\left\{\forall i,\ \norm{\mu_{P_i}-\mu_{\widehat P_i}}_{\Hcal}\le \tau_i(\delta/N)\right\}\ge 1-\delta.
\end{equation}
Thus all true embeddings lie in their Hilbert ambiguity sets simultaneously with probability at least $1-\delta$.

\subsection{Empirical CVaR Used in the Code}

For active-learning losses $e_1,\ldots,e_T$, sort them in descending order as $e_{(1)}\ge\cdots\ge e_{(T)}$. The empirical CVaR at tail fraction $\alpha$ is
\begin{equation}
    \widehat{\CVaR}_\alpha=\frac{1}{m}\sum_{r=1}^{m}e_{(r)},\qquad m=\lceil\alpha T\rceil.
\end{equation}
This is the quantity reported in the supplied CSV files.

\section{Submission Checklist}
\label{app:checklist}

For a final submission:
\begin{itemize}[leftmargin=1.1em]
    \item keep the Heart-of-Risk score as the main object, not as a motivational preface;
    \item state clearly that KME is a conservative envelope of transport ambiguity, not an exact replacement;
    \item add the exact PMARLDB/PMARLDS thesis solvers to the supplied benchmark interface;
    \item run 50 repeats for GAUSSDIST, digit distributions, and 20 Newsgroups;
    \item report AULC, final accuracy, empirical CVaR, worst-round loss, average queried radius, runtime, and average rank.
\end{itemize}

\balance
\bibliographystyle{IEEEtran}
\begin{thebibliography}{99}

\bibitem{settles2012active}
B. Settles, \emph{Active Learning}. Morgan \& Claypool, 2012.

\bibitem{ghafarian2019fg}
S. H. Ghafarian and H. Sadoghi Yazdi, ``Functional gradient approach to probabilistic minimax active learning,'' \emph{Engineering Applications of Artificial Intelligence}, vol. 85, pp. 21--32, 2019.

\bibitem{ghafarian2022cybernetics}
S. H. Ghafarian and H. Sadoghi Yazdi, ``Prepare for the worst, hope for the best: Active robust learning on distributions,'' \emph{IEEE Transactions on Cybernetics}, vol. 52, no. 6, pp. 5573--5586, 2022.

\bibitem{smola2007hilbert}
A. Smola, A. Gretton, L. Song, and B. Sch{\"o}lkopf, ``A Hilbert space embedding for distributions,'' in \emph{Algorithmic Learning Theory}, 2007, pp. 13--31.

\bibitem{gretton2012mmd}
A. Gretton, K. M. Borgwardt, M. J. Rasch, B. Sch{\"o}lkopf, and A. Smola, ``A kernel two-sample test,'' \emph{Journal of Machine Learning Research}, vol. 13, pp. 723--773, 2012.

\bibitem{muandet2017kme}
K. Muandet, K. Fukumizu, B. Sriperumbudur, and B. Sch{\"o}lkopf, ``Kernel mean embedding of distributions: A review and beyond,'' \emph{Foundations and Trends in Machine Learning}, vol. 10, no. 1--2, pp. 1--141, 2017.

\bibitem{sriperumbudur2010hilbert}
B. K. Sriperumbudur, A. Gretton, K. Fukumizu, G. Lanckriet, and B. Sch{\"o}lkopf, ``Hilbert space embeddings and metrics on probability measures,'' \emph{Journal of Machine Learning Research}, vol. 11, pp. 1517--1561, 2010.

\bibitem{szabo2016distribution}
Z. Szab{\'o}, A. Gretton, B. P{\'o}czos, and B. Sriperumbudur, ``Two-stage sampled learning theory on distributions,'' in \emph{Artificial Intelligence and Statistics}, 2015, pp. 948--957.

\bibitem{flaxman2016bayesian}
S. Flaxman, D. Sejdinovic, J. P. Cunningham, and S. Filippi, ``Bayesian learning of kernel embeddings,'' in \emph{Uncertainty in Artificial Intelligence}, 2016.

\bibitem{villani2009ot}
C. Villani, \emph{Optimal Transport: Old and New}. Springer, 2009.

\bibitem{mohajerin2018wasserstein}
P. Mohajerin Esfahani and D. Kuhn, ``Data-driven distributionally robust optimization using the Wasserstein metric: Performance guarantees and tractable reformulations,'' \emph{Mathematical Programming}, vol. 171, pp. 115--166, 2018.

\bibitem{kuhn2019wasserstein}
D. Kuhn, P. Mohajerin Esfahani, V. A. Nguyen, and S. Shafieezadeh-Abadeh, ``Wasserstein distributionally robust optimization: Theory and applications in machine learning,'' arXiv:1908.08729, 2019.

\bibitem{staib2019mmd}
M. Staib and S. Jegelka, ``Distributionally robust optimization and generalization in kernel methods,'' in \emph{Advances in Neural Information Processing Systems}, 2019.

\bibitem{feydy2019sinkhorn}
J. Feydy, T. S{\'e}journ{\'e}, F.-X. Vialard, S.-i. Amari, A. Trouv{\'e}, and G. Peyr{\'e}, ``Interpolating between optimal transport and MMD using Sinkhorn divergences,'' in \emph{Artificial Intelligence and Statistics}, 2019, pp. 2681--2690.

\bibitem{genevay2019sample}
A. Genevay, L. Chizat, F. Bach, M. Cuturi, and G. Peyr{\'e}, ``Sample complexity of Sinkhorn divergences,'' in \emph{Artificial Intelligence and Statistics}, 2019, pp. 1574--1583.

\bibitem{vayer2023controlling}
T. Vayer and R. Gribonval, ``Controlling Wasserstein distances by kernel norms with application to compressive statistical learning,'' \emph{Journal of Machine Learning Research}, vol. 24, no. 149, pp. 1--51, 2023.

\bibitem{sriperumbudur2009kernel}
B. K. Sriperumbudur, K. Fukumizu, A. Gretton, G. Lanckriet, and B. Sch{\"o}lkopf, ``Kernel choice and classifiability for RKHS embeddings of probability distributions,'' in \emph{Advances in Neural Information Processing Systems}, 2009.

\bibitem{huang2014active}
S.-J. Huang, R. Jin, and Z.-H. Zhou, ``Active learning by querying informative and representative examples,'' \emph{IEEE Transactions on Pattern Analysis and Machine Intelligence}, vol. 36, no. 10, pp. 1936--1949, 2014.

\bibitem{zhu2003semi}
X. Zhu, Z. Ghahramani, and J. Lafferty, ``Semi-supervised learning using Gaussian fields and harmonic functions,'' in \emph{International Conference on Machine Learning}, 2003.

\bibitem{scheffer2001active}
T. Scheffer, C. Decomain, and S. Wrobel, ``Active hidden Markov models for information extraction,'' in \emph{International Symposium on Intelligent Data Analysis}, 2001.

\bibitem{chang2011libsvm}
C.-C. Chang and C.-J. Lin, ``LIBSVM: A library for support vector machines,'' \emph{ACM Transactions on Intelligent Systems and Technology}, vol. 2, no. 3, pp. 1--27, 2011.

\bibitem{mikolov2013distributed}
T. Mikolov, K. Chen, G. Corrado, and J. Dean, ``Efficient estimation of word representations in vector space,'' arXiv:1301.3781, 2013.


\bibitem{artzner1999coherent}
P. Artzner, F. Delbaen, J.-M. Eber, and D. Heath, ``Coherent measures of risk,'' \emph{Mathematical Finance}, vol. 9, no. 3, pp. 203--228, 1999.

\bibitem{rockafellar2000cvar}
R. T. Rockafellar and S. Uryasev, ``Optimization of conditional value-at-risk,'' \emph{Journal of Risk}, vol. 2, pp. 21--41, 2000.

\end{thebibliography}

\end{document}
